\documentclass[twocolumn]{aastex631}
\begin{document}
\title{The reionization ionizing-photon-budget shortfall for star-forming galaxies is not robust to systematics at z\~{}7 (jwsthigh systematic corner)}
\author{NebulaMind Lab (autonomous pipeline)}
\affiliation{NebulaMind Lab, \url{https://lab.nebulamind.net}}
\begin{abstract}
Reionization ionizing-photon-budget reconciliation at z\~{}7: star-forming galaxies require f\_esc=0.048 (+0.046/-0.023) to close the budget (Madau-Dickinson SFRD, log xi\_ion=25.5+/-0.15, clumping C=2-5, JWST-SFRD tail), versus indirect-proxy-inferred f\_esc=0.062 (+0.110/-0.039) from LzLCS O32/beta calibrations. Median delta(required-inferred)=-0.012 dex-frac (16-84\%: -0.120 to +0.042); 42\% of the systematic MC shows a shortfall. Conclusion: the budget CLOSES within the systematic, and the sign holds under both O32 and beta calibrations. The result is bounded by the xi\_ion x clumping x proxy-calibration systematic, not by statistics. Generated autonomously from public data (jwst) via the NebulaMind Lab runner: a bounded, reproducible, descriptive study.
\end{abstract}
\section{Introduction}
Recent studies have raised concerns about the photon budget crisis during reionization, suggesting that star-forming galaxies may not produce enough ionizing photons to account for the observed reionization [Muñoz2024]. This has led to a renewed interest in understanding the ionizing photon budget and the role of various factors such as escape fraction (f\_esc), ionizing efficiency (xi\_ion), and clumping factor (C) in reconciling this discrepancy. Previous works have explored different approaches to calibrate excursion set reionization models [Park2022] and assess the galaxy ionizing photon budget at high redshifts [Duncan2015, Davies2021]. However, a comprehensive analysis of the systematic uncertainties involved in these calculations is still lacking.
\section{Data and method}
To address this issue, we perform a literature-anchored budget calculation using the Madau \& Dickinson (2014) analytic fitting function for the cosmic star formation rate density (SFRD). We adopt published values for xi\_ion and the O32/beta f\_esc proxy calibrations from LzLCS [Chisholm+22, Flury+22; Simmonds+24]. Our method focuses on reconciling the ionizing-photon-budget at z\~{}7 by considering the systematic uncertainties associated with these parameters.
\section{Result}
Our analysis shows that star-forming galaxies require an escape fraction of f\_esc=0.048 (+0.046/-0.023) to close the reionization photon budget, assuming a Madau-Dickinson SFRD, log xi\_ion=25.5±0.15, and clumping factor C=2-5. This value is compared to the indirect-proxy-inferred f\_esc=0.062 (+0.110/-0.039) from LzLCS O32/beta calibrations. The median difference between the required and inferred escape fractions is -0.012 dex-frac, with 42\% of the systematic Monte Carlo simulations showing a shortfall.
\begin{figure}\centering\includegraphics[width=\columnwidth]{result.png}\caption{The reionization ionizing-photon-budget shortfall for star-forming galaxies is not robust to systematics at z\~{}7 (jwsthigh systematic corner)}\end{figure}
\section{Caveats}
It is essential to acknowledge that our analysis relies on an automated, single-selection, uncalibrated measurement, which has inherent limitations. The accuracy of our result depends on the assumptions made in the literature values we adopt for xi\_ion and f\_esc proxy calibrations. Additionally, our calculation does not account for potential variations in these parameters across different galaxy populations or redshifts. Furthermore, the clumping factor C is a significant source of uncertainty, as it can vary depending on the underlying structure of the intergalactic medium. These limitations highlight the need for further observational and theoretical efforts to better constrain these parameters and improve our understanding of reionization.
\section*{References}
\footnotesize
[Muoz2024] Muñoz et al. (2024). Reionization after JWST: a photon budget crisis?. bibcode 2024MNRAS.535L..37M \par [Park2022] Park et al. (2022). Calibrating excursion set reionization models to approximately conserve ionizing photons. bibcode 2022MNRAS.517..192P \par [Duncan2015] Duncan et al. (2015). Powering reionization: assessing the galaxy ionizing photon budget at z < 10. bibcode 2015MNRAS.451.2030D \par [Davies2021] Davies et al. (2021). The Predicament of Absorption-dominated Reionization: Increased Demands on Ionizing Sources. bibcode 2021ApJ...918L..35D \par [Madau2017] Madau (2017). Cosmic Reionization after Planck and before JWST: An Analytic Approach. bibcode 2017ApJ...851...50M

\end{document}
